\section*{Chapter \ref{ch:ll:rust-for-low-latency} --- Rust for Low Latency}

\begin{solution}{exo:ll:rust-for-low-latency:1}
On the committed measurement about 10.3 and \qty{10.6}{\nano\second}: Rust within a few percent of C++ (a fifth at most
between runs). At 5 million messages a second a difference of \qty{0.3}{\nano\second} is 1.5 milliseconds a second, noise
next to the path's other costs.
\end{solution}

\begin{solution}{exo:ll:rust-for-low-latency:2}
\texttt{\textquotedbl{}ESZ6\textquotedbl{}.to\_string()} allocates (a \texttt{String} has no inline buffer); \texttt{Vec::new()} does not; its first
\texttt{push} does (room for four small elements); \texttt{[0u8; 24]} is on the stack; \texttt{Box::new(3)} allocates.
\end{solution}

\begin{solution}{exo:ll:rust-for-low-latency:3}
Rust: the symbol's \texttt{String} (1, where C++'s inline buffer avoided it), the identifier's \texttt{String} (1, as in
C++), the vector (1: the first push reserves four, against three growth steps in C++), the map entry (0: the key is
moved in and the table was already allocated, against a node and a key copy in C++), the clone of the fills (1, as C++'s
lambda copy), the boxed closure (1, against two in \texttt{std::function}). $1+1+1+0+1+1 = 5$ against $0+1+3+2+1+2 = 9$.
\end{solution}

\begin{solution}{exo:ll:rust-for-low-latency:4}
Four: two 64-bit loads of two \texttt{u32} each, widened to 64 bits and added into two vector accumulators of two lanes.
The remainder of zero to three elements goes through the scalar loop at \texttt{.LBB16\_4}.
\end{solution}

\begin{solution}{exo:ll:rust-for-low-latency:5}
Each iteration waits for a load at a scattered address; the check is an independent comparison with a predicted branch,
executed in the shadow of that wait. It would show in a loop bound by the processor rather than by memory, over data in the
first-level cache, and above all where the check prevents vectorisation.
\end{solution}

\begin{solution}{exo:ll:rust-for-low-latency:6}
No. \texttt{CheckedIndex::new(\&[99], 100)} followed by \texttt{sum(\&v, 100)} with \texttt{v} of length 10 would read
\texttt{v[99]} out of bounds, in safe code: the abstraction would have trusted its caller about \texttt{len}. Comparing the
slice's actual length is what makes it \omterm{def:ll:rust-for-low-latency:sound}{sound}.
\end{solution}

\begin{solution}{exo:ll:rust-for-low-latency:7}
The explicit \texttt{match} makes the error case visible: \texttt{flatten()} silently drops a decode error, while the
\texttt{match} can count it and stop. Neither version allocates (the decoder borrows the buffer), and neither contains a
panic path in the loop body when the lengths are checked by the decoder's own comparisons; the difference is behaviour on
bad input, not speed.
\end{solution}

\begin{solution}{exo:ll:rust-for-low-latency:8}
A market-data handler is a single ordered stream: it cannot be spread over cores without losing order, and a work-stealing
scheduler adds queueing, migration between threads and cache misses, exactly the \omterm{def:ll:where-latency-comes-from:tail}{jitter} the hot path avoids. Use a pinned
thread spinning on its socket or ring; keep async for the many slow connections.
\end{solution}

\begin{solution}{pb:ll:rust-for-low-latency:1}
\begin{enumerate}
\item About \qty{10.3}{\nano\second} (C++) and \qty{10.6}{\nano\second} (Rust) on the committed measurement: a ratio of about
1.03, and within a fifth run to run.
\item LLVM vectorises the counted loop at \texttt{-O}; g++ 11 does not vectorise at \texttt{-O2} (chapter~14).
\item \texttt{-O3}, or g++ 12 or later at \texttt{-O2}, or an explicitly vectorised loop.
\item That they cost nothing measurable there: checked and unchecked gathers take the same time within noise in both
languages.
\item Five (Rust) and nine (C++) for the drafts; none for either fixed handler.
\item \texttt{String}, which always allocates, even for a short symbol.
\item A counting \omterm{def:ll:rust-for-low-latency:alloc}{global allocator} installed in the test binary, and an assertion on the count per message after warm-up.
\item The count rises and the test fails, provided the test drives the dependency's hot paths with production-shaped data.
\item \omterm{def:ll:rust-for-low-latency:race}{Data races}, use-after-free and dangling references, and out-of-bounds reads (turned into panics).
\item In calls to the C library (affinity), in lock-free structures (chapter~11) and in unchecked indexing where a \omterm{def:ll:rust-for-low-latency:sound}{sound
abstraction} justifies it.
\item That no sequence of calls through its safe interface can cause \omterm{def:ll:cpp-for-latency-iii-the-compiler:ub}{undefined behaviour}.
\item A bug has been caught before corrupting memory; the process should stop quoting (cancel orders, chapter~22) and
restart, not continue.
\item \textbf{Named result.} The Rust decoder runs at about 1.03 times the C++ decoder's time per message on the committed
measurement (within a fifth run to run); the drafts allocate five (Rust) and nine (C++) times per message, the fixed
handlers none.
\item No: the difference is within run-to-run noise and small against the path's other stages.
\item The feed decoder and book builder: stateful parsing of untrusted input, where memory errors are likeliest and most
costly.
\item Around the hot path: venue REST and websocket connections, drop copies, internal services.
\item Through the plan file written by \texttt{firm.coreplan} and read by both \texttt{firm.affinity} twins, and through shared
fixtures that both implementations must reproduce.
\item The existing C++ code, tooling and experience; some libraries; and time spent on the port instead of features.
\item Run both on the same recorded inputs and compare outputs message by message (differential testing, chapter~25).
\item It moves memory and concurrency errors from production to the compiler, at the same speed.
\end{enumerate}
\end{solution}

\begin{solution}{iq:ll:rust-for-low-latency:1}
The abstraction compiles to what one would write by hand. Verify by reading the optimised assembly of the chain and of the
hand-written loop, and by benchmarking both on realistic data.
\iqlookfor{assembly and measurement, not faith.}
\end{solution}

\begin{solution}{iq:ll:rust-for-low-latency:2}
Rarely: a predicted check is cheap, and the compiler removes checks it can prove. They cost when they prevent
vectorisation or sit in a processor-bound loop. Remove them safely with iterators, with an assertion on the length before
the loop, or behind a \omterm{def:ll:rust-for-low-latency:sound}{sound abstraction} that validates indices once.
\iqlookfor{elimination by proof, and \omterm{def:ll:rust-for-low-latency:sound}{sound} unchecked access.}
\end{solution}

\begin{solution}{iq:ll:rust-for-low-latency:3}
Its safe interface cannot be used to cause \omterm{def:ll:cpp-for-latency-iii-the-compiler:ub}{undefined behaviour}: every invariant the unsafe code needs is established and
checked by the abstraction itself, not assumed of callers; each unsafe block states its argument.
\iqlookfor{invariants owned by the abstraction.}
\end{solution}

\begin{solution}{iq:ll:rust-for-low-latency:4}
Mutable references are exclusive and only \texttt{Send}/\texttt{Sync} types cross threads, checked at compile time. It does
not prevent race conditions in logic (two atomic operations that should have been one), deadlocks, or races inside
\texttt{unsafe} code.
\iqlookfor{\omterm{def:ll:rust-for-low-latency:race}{data races} versus race conditions.}
\end{solution}

\begin{solution}{iq:ll:rust-for-low-latency:5}
No: a runtime schedules tasks across threads and interleaves them, adding latency and \omterm{def:ll:where-latency-comes-from:tail}{jitter}; the hot path is a pinned
thread spinning on its input. Async fits the many waiting connections around it.
\iqlookfor{scheduling \omterm{def:ll:where-latency-comes-from:tail}{jitter} versus a dedicated thread.}
\end{solution}

\begin{solution}{iq:ll:rust-for-low-latency:6}
Install a counting \omterm{def:ll:rust-for-low-latency:alloc}{global allocator} in a test binary (it sees every allocation, including those of dependencies), drive the
component with production-shaped inputs beyond warm-up, and assert that the count does not change; run it in continuous
integration.
\iqlookfor{the \omterm{def:ll:rust-for-low-latency:alloc}{global allocator} as the single choke point.}
\end{solution}
